% =====================================================================
%  LỜI NÓI ĐẦU
% =====================================================================
\chapter*{LỜI NÓI ĐẦU}
\addcontentsline{toc}{chapter}{LỜI NÓI ĐẦU}

Mua sắm, bảo quản thực phẩm và chuẩn bị bữa ăn là những công việc thường nhật của mỗi gia đình. Khi các hoạt động này được quản lý bằng trí nhớ hoặc ghi chú rời rạc, người dùng có thể bỏ sót mặt hàng, mua trùng hoặc để thực phẩm quá hạn. Theo Chương trình Môi trường Liên Hợp Quốc, hộ gia đình chiếm khoảng 60\% lượng thực phẩm bị lãng phí trên toàn thế giới~\cite{unep2024}. Từ thực tế đó, nhóm chọn đề tài xây dựng ứng dụng di động đa nền tảng “Đi chợ tiện lợi” cho học phần Phát triển ứng dụng đa nền tảng.

Ứng dụng “Đi chợ tiện lợi” kết nối danh sách mua sắm, tủ lạnh và kế hoạch bữa ăn. Người dùng có thể lập danh sách theo ngày, chia sẻ trong gia đình và nhận phân công từ trưởng nhóm. Thực phẩm đã mua được ghi nhận vào tủ cùng hạn dùng đề xuất; hệ thống gửi nhắc nhở theo ngưỡng do người dùng đặt (mặc định ba ngày). Kế hoạch bữa ăn dựa trên tồn kho, ưu tiên thực phẩm cần dùng sớm. Các báo cáo tổng hợp thói quen mua sắm, tiêu thụ và giá trị thực phẩm bị lãng phí. Ứng dụng hướng tới hoạt động trên Android và iOS từ cùng một mã nguồn.

Nhóm thực hiện đề tài qua các bước: khảo sát thực trạng và giải pháp hiện có, phân tích nghiệp vụ, đặc tả yêu cầu bằng use case theo UML, thiết kế, xây dựng và kiểm thử ứng dụng. Các chức năng mở rộng gồm tạo danh sách mua sắm từ kế hoạch bữa ăn, tra cứu mã vạch, quản lý tồn kho theo nguyên tắc hết hạn trước dùng trước (FEFO) và thống kê lãng phí thực phẩm.

Ở giai đoạn hiện tại, báo cáo gồm hai chương đã hoàn thiện. Chương 1 khảo sát bài toán và phân tích nghiệp vụ; Chương 2 đặc tả tác nhân, use case và yêu cầu phi chức năng. Phụ lục gồm từ điển dữ liệu, đặc tả use case bổ sung, danh mục thông điệp, ma trận truy vết và mã nguồn PlantUML.

Nhóm xin cảm ơn thầy Nguyễn Quốc Tuấn đã hướng dẫn trong quá trình học tập và thực hiện bài tập lớn. Do thời gian và kinh nghiệm còn hạn chế, báo cáo có thể còn thiếu sót. Nhóm mong nhận được nhận xét và góp ý của thầy để hoàn thiện đề tài.
